% page 415 du fac-simile (image p-414 du scan)
% bloc 157.5 x 246.3 mm ; marge gauche 44.5 mm
\pgc{22.860mm}{0.000mm}{
\l{\cel{50}407}
\l{}
\l{\sou{oosporo}. (\sou{bot.})\cel{2}Nomo donita ulakaze a la ovo di la algi o}
\l{\cel{3}di la fungi, precipue kande la organo feminala (oosfero)}
\l{\cel{3}liberigesas ante fekundigeso, nam lore la ovo kelke si-}
\l{\cel{3}milesas sporo. - DEFIRS.}
\l{}
\l{\sou{-op-}\cel{1}. Sufixo distributala, kun nomi di nombro o di quanto.}
\l{}
\l{\sou{opaka}. Qua ne lasas la lumo-radii pasar tra su. - DEFIS.}
\l{}
\l{\sou{opalo}. (\sou{geol.}) Lapido volkanala, lakto-blanka e bluatra,}
\l{\cel{3}kun reflekti koloro-chanjanta. - DEFIRS.}
\l{}
\l{\sou{opalecar}\cel{2}(\sou{netrans.})\cel{2}Qua aquiras nuanco opalea. DEFIS.}
\l{}
\l{\sou{opero}. Verko teatrala e muzikala, qua konsistas ye recitativi,}
\l{\cel{3}kanti helpata da orkestro, ed ulakaze kun dansi. - DEFIRS.}
\l{}
\l{\sou{operacar}. I. (\sou{trans.}) (\sou{kirurgio}) Submisar ad interveno kirur-}
\l{\cel{3}giala. - II. (\sou{netrans.)}\cel{2}(1) (\sou{matem.}) Agar la kalkuli quin}
\l{\cel{3}onu efektigas departante de quanti konocita, por renkontrar}
\l{\cel{3}un o plura quanti nekonocita. (2) (\sou{militarto}) A g a r}
\l{\cel{3}la trupo-diplasi qui havas kom skopo sucesigar atako o de-}
\l{\cel{3}fenso. - (3) (en\cel{1}\sou{borso financala o komercala}) Manovrar por}
\l{\cel{3}sucesigar profitoze la kompro o la vendo di la komercajo.}
\l{\cel{3}- DEFIRS.}
\l{}
\l{\sou{operkulo}. (\sou{tekn.)}\cel{3}Aparato destinita ad obturar ula orifici.}
\l{\cel{3}EFIS.}
\l{}
\l{\sou{opiato}.\cel{2}I. Elektuario opiumoza. - II. Irga elektuario, pasto}
\l{\cel{3}ek pudro dilutita en siropo. - DEFIS.}
\l{}
\l{\sou{opinionar}. (\sou{netrans.})\cel{2}Enuncar to quon onu judikas (sen funda-}
\l{\cel{3}mento ciencale establisita) pri temo o persono. - EFIS.}
\l{}
\l{\sou{opiumo.}\cel{2}Substanco narkotigiva, kalmigiva se uzata moderata-}
\l{\cel{3}doze, ma venena granda-doze, quan onu extraktas de la suko}
\l{\cel{3}di ula sorti de papaveri. - DEFIRS.}
\l{}
\l{-opl-.\cel{2}Sufixo pri multipliko.}
\l{}
\l{opodeldoko. (\sou{farmacio}) Medikamento qua havas kom fundamento}
\l{\cel{3}solvuro alkoholoza ek sapono animalala o medicinala, a qua}
\l{\cel{3}onu adjuntas kamforo, amoniako ed esenco de timiano e de}
\l{\cel{3}rosmarino,uzata fricione(kontre la dolori). - DEFI?}
\l{}
\l{oportuna. Qua eventas precize kande lo esas bona por to quo}
\l{\cel{3}koncernesas. - DEFIS.}
\l{}
\l{oportar. (\sou{netrans.}) Esar absolute necesa. - L.}
\l{}
\l{opozar.(trans.) I. Lokizar (ulo) regarde ulo. - II. Lokizar}
\l{\cel{3}(ulu, ulo) regarde altra, por esar obstaklo pri ica.}
\l{\cel{3}- III. Alegar (ulo) kontre to quon dicas o pensas altru.}
\l{\cel{55}DEFIRS.}
}
